\documentclass[11pt]{article}
\usepackage[margin=1in]{geometry}
\usepackage{times}
\usepackage{latexsym}
\usepackage{amsmath}
\usepackage{amssymb}
\usepackage{booktabs}
\usepackage{multirow}
\usepackage{graphicx}
\usepackage{xcolor}
\usepackage{url}
\usepackage[colorlinks=true,citecolor=blue,urlcolor=blue]{hyperref}
\usepackage{natbib}
\renewcommand{\bibsection}{\section*{References}}

\title{How Much Contamination Are We Actually Correcting? \\
A Joint Psychometric Audit of LLM-as-a-Judge Evaluation}

\author{Anonymous}

\begin{document}
\maketitle

\begin{abstract}
LLM-as-a-judge evaluation faces two entangled threats: judge-specific scoring bias and benchmark contamination through training data overlap.
Existing methods address these independently---bias-only corrections (e.g., CyclicJudge, UDA) ignore contamination, while contamination-only adjustments (e.g., DCR) ignore judge variation---yet the two effects confound each other when estimated sequentially.
We introduce the \textbf{Gated Many-Facet Rasch Model (G-MFRM)}, a staged joint calibration framework that simultaneously estimates model ability ($\theta_n$), item difficulty ($b_i$), judge severity ($d_j$), and contamination susceptibility ($c_n$) within a unified psychometric model.
We formalize the confounding between judge severity and contamination as an omitted variable bias (OVB) problem and show via Monte Carlo simulation that the resulting $\theta$-bias scales linearly with contamination magnitude.
Applying G-MFRM to 8 LLM judges $\times$ 22 target models across MT-Bench, Arena-Hard, and WildBench, we find a striking null result: real contamination on current benchmarks is negligible (median $c_n = 0.15$, far below synthetic assumptions of $c_n \geq 0.5$), and no method---including G-MFRM---significantly outperforms standard multi-facet calibration for ranking (all $\Delta\tau < 0.02$, $p > 0.18$).
We argue this null finding is itself a contribution: it challenges widespread contamination assumptions, validates G-MFRM as a diagnostic framework that \emph{confirms} benchmark cleanliness, and establishes contamination magnitude thresholds ($c_n \geq 0.7$) above which joint correction becomes necessary.
Through adversarial injection experiments, we demonstrate G-MFRM achieves $\tau = 0.87$ versus $0.75$ for uncorrected baselines when contamination is artificially elevated, confirming the method works when contamination exists.
\end{abstract}

\section{Introduction}

LLM-as-a-judge has become the dominant paradigm for scalable model evaluation \cite{zheng2023judging}, yet two systematic threats undermine the validity of judge-based rankings.
First, judges exhibit scoring biases---position bias, verbosity bias, and idiosyncratic severity \cite{wang2023large}---that distort relative model rankings.
Second, benchmark contamination through training data overlap inflates scores on memorized items \cite{golchin2023time,shi2023detecting}, creating a confound between genuine ability and memorization.

These threats have been studied in separate literatures.
Bias-correction methods like CyclicJudge \cite{li2024cyclicjudge} and UDA \cite{huang2024uda} eliminate inter-judge variance but are contamination-unaware.
Contamination detection methods like DCR \cite{yang2023dcr} and ConStat \cite{dekoninck2024constat} quantify memorization effects but ignore judge variation.
A sequential pipeline---first correct bias, then adjust for contamination (or vice versa)---appears natural but introduces systematic error: the first stage's estimates are biased by the unmodeled facet, and this bias propagates to the second stage.

We formalize this confounding as an omitted variable bias (OVB) problem (Section~\ref{sec:ovb}).
When judge severity $d_j$ and contamination susceptibility $c_n \cdot e_i$ are correlated, omitting either from the model produces asymptotically biased estimates of the other and, consequently, of model ability $\theta_n$.
Via Monte Carlo simulation across 48 conditions, we show the bias magnitude scales linearly with contamination: at $c_n = 1.5$, the $\theta$-RMSE degradation reaches 43.8\% compared to joint estimation, while at the levels observed in real data ($c_n \approx 0.15$), bias is negligible at 0.02 logits (Section~\ref{sec:confounding}).

To resolve this confounding, we propose the \textbf{Gated Many-Facet Rasch Model (G-MFRM)}, a three-stage framework that extends MFRM \cite{linacre1989many,eckes2011introduction} with a contamination gating mechanism inspired by DG-IRT \cite{shu2013item}.
G-MFRM estimates all facets---ability $\theta_n$, difficulty $b_i$, severity $d_j$, and contamination susceptibility $c_n$---in a staged joint calibration, producing bias-and-contamination-corrected ``fair scores'' alongside interpretable diagnostic outputs.

Applying G-MFRM to 8 LLM judges $\times$ 22 target models across three benchmarks (MT-Bench, Arena-Hard, WildBench), we obtain a set of results that, while initially disappointing for the ranking-improvement narrative, constitute three substantive contributions:

\begin{enumerate}
\item \textbf{A calibrated null finding}: Real contamination on current benchmarks is negligible (median $c_n = 0.15$), far below levels assumed in contamination-correction literature ($c_n \geq 0.5$). No method significantly outperforms standard MFRM calibration ($\Delta\tau < 0.02$, all $p > 0.18$). This challenges implicit assumptions that training data overlap produces large performance inflation on modern, regularly-refreshed benchmarks.

\item \textbf{A principled diagnostic framework}: G-MFRM provides four-dimensional output---per-model $c_n$, per-item $e_i$, per-judge $d_j$, and corrected $\theta_n$---that no existing tool combines. Even when ranking correction is unnecessary, these diagnostics enable benchmark auditing, judge calibration assessment, and contamination monitoring.

\item \textbf{Methodology with demonstrated validity}: Through adversarial injection experiments, we show G-MFRM achieves $\tau = 0.87$ versus $0.75$ for uncorrected baselines at elevated contamination ($c_n = 1.5$), with a crossover point at $c_n \approx 0.7$. The method works when contamination exists; it simply does not exist at actionable levels on current benchmarks.
\end{enumerate}

\section{Related Work}

\paragraph{LLM judge bias correction.}
Zheng et al.~\cite{zheng2023judging} introduced MT-Bench and documented position bias in LLM judges.
CyclicJudge \cite{li2024cyclicjudge} eliminates average judge bias through round-robin evaluation.
UDA \cite{huang2024uda} applies neural K-factor adjustment, reducing inter-judge standard deviation by 63.4\%.
Preference Leakage \cite{ryu2024prefleakage} quantifies same-family bias (PLS $\approx$ 23.6\%).
These methods correct bias effectively but assume contamination is absent or orthogonal to bias.

\paragraph{Benchmark contamination detection.}
Golchin \& Surdeanu \cite{golchin2023time} and Shi et al.~\cite{shi2023detecting} demonstrate that training data overlap inflates benchmark scores.
DCR \cite{yang2023dcr} factors out contamination effects ($r^2 = 0.94$).
ConStat \cite{dekoninck2024constat} provides per-model contamination statistics.
These approaches quantify contamination but do not account for judge variation.

\paragraph{Psychometric models for evaluation.}
Many-Facet Rasch Measurement (MFRM) \cite{linacre1989many,eckes2011introduction,engelhard2013invariant} models examinee ability, item difficulty, and rater severity within a single calibration.
In educational testing, joint modeling of rater effects and differential item functioning (DIF) is established \cite{wind2018exploring,patz2002testlet,fox2001bayesian}.
DG-IRT \cite{shu2013item} introduces gating mechanisms for item compromise detection.
G-MFRM adapts these psychometric traditions to LLM evaluation, replacing binary compromise indicators with continuous exposure estimates.

\section{Method}

\subsection{Problem Formulation}

Let $X_{nij}$ denote the score assigned by judge $j$ to model $n$ on item $i$.
We seek to estimate model ability $\theta_n$ free from both judge bias $d_j$ and contamination inflation $c_n \cdot e_i$, where $e_i \in [0,1]$ indicates item exposure and $c_n \geq 0$ captures model-level contamination susceptibility.

\subsection{Omitted Variable Bias in Sequential Estimation}
\label{sec:ovb}

Consider the data-generating process:
\begin{equation}
X_{nij} = \theta_n - b_i + d_j + c_n \cdot e_i + \epsilon_{nij}
\end{equation}
where $\epsilon_{nij} \sim \mathcal{N}(0, \sigma^2)$.
A bias-only model omits $c_n \cdot e_i$; a contamination-only model omits $d_j$.
Following standard OVB analysis \cite{angrist2009mostly}, when $\text{Cov}(d_j, c_n \cdot e_i | \theta_n, b_i) \neq 0$, the bias-only estimator $\hat{\theta}_n^{\text{bias-only}}$ satisfies:
\begin{equation}
\text{plim}\,\hat{\theta}_n^{\text{bias-only}} = \theta_n + \gamma \cdot c_n \cdot \bar{e}
\end{equation}
where $\gamma$ captures the non-orthogonality between judge and contamination facets.
The direction of bias is determined by $\text{sign}(\gamma)$: lenient judges evaluating contamination-inflated models produce upward-biased $\hat{\theta}_n$.
Sequential pipelines do not resolve this because the first stage's biased estimates serve as inputs to the second.

\subsection{Gated Many-Facet Rasch Model}

G-MFRM proceeds in three stages:

\paragraph{Stage 1: Standard MFRM.}
Fit a standard MFRM to obtain initial estimates $\hat{\theta}_n^{(0)}$, $\hat{b}_i$, $\hat{d}_j$:
\begin{equation}
\log \frac{P(X_{nij} \geq k)}{P(X_{nij} < k)} = \theta_n - b_i + d_j - \tau_k
\end{equation}
where $\tau_k$ are threshold parameters. This stage is well-identified and provides calibrated judge severity estimates.

\paragraph{Stage 2: External exposure estimation.}
Estimate per-item exposure $e_i$ via a three-source pipeline:
(a) \emph{Performance gap}: difference between scores on original items vs.\ reference variants (rephrased items preserving difficulty), with spline correction for difficulty confounds;
(b) \emph{Temporal signal}: score changes across model release dates;
(c) \emph{LogProber}: direct memorization probing where logprob APIs are available.
Sources are combined via weighted composite. This stage uses external information not present in $X_{nij}$, ensuring $e_i$ is not confounded with the scoring model.

\paragraph{Stage 3: Joint gated model.}
With $e_i$ fixed from Stage 2, fit the full model:
\begin{equation}
\log \frac{P(X_{nij} \geq k)}{P(X_{nij} < k)} = \theta_n - b_i + d_j + c_n \cdot e_i - \tau_k
\label{eq:gmfrm}
\end{equation}
with priors $\theta_n \sim \mathcal{N}(0, \sigma_\theta^2)$, $c_n \sim \text{Half-Cauchy}(0, 1)$, $d_j \sim \mathcal{N}(0, \sigma_d^2)$.
The ``fair score'' is $\hat{\theta}_n^{\text{fair}} = \hat{\theta}_n^{(3)}$, estimated jointly with all facets.
The staged approach---fixing $e_i$ rather than estimating it jointly---is a deliberate identifiability choice analogous to two-step estimators in econometrics \cite{heckman1979sample}.

\paragraph{Identifiability.}
With $e_i$ fixed, the model parameters $(\theta_n, b_i, d_j, c_n)$ are identifiable under standard conditions: the design matrix has full column rank when there exist items with $e_i > 0$ and $e_i \approx 0$, and when multiple judges evaluate each model-item combination.

\paragraph{Diagnostic outputs.}
Beyond corrected rankings, G-MFRM produces: (a) per-model $c_n$ for contamination auditing, (b) per-item $e_i$ for item retirement decisions, (c) per-judge $d_j$ for judge calibration monitoring, and (d) person-fit statistics (infit/outfit) for detecting aberrant response patterns.

\section{Experiments}

\subsection{Data Collection}
\label{sec:data}

We collect multi-judge evaluation data using 8 LLM judges (GPT-5.5, Claude Opus 4.7, Claude Sonnet 4.6, DeepSeek-V4-Pro, Qwen3-235B-A22B, Qwen3.6-27B, Llama 4 Maverick, Gemini 2.5 Pro) evaluating 22 target models on MT-Bench (80 items, 14{,}080 judgments) and Arena-Hard (500 items, 35{,}200 judgments).
Position-swapped pairs yield agreement of 0.72.
Inter-judge agreement is moderate (Fleiss' $\kappa = 0.38$), reflecting genuine disagreement on style and difficulty rather than contamination-driven divergence.
Table~\ref{tab:judges} summarizes judge characteristics.

\begin{table}[t]
\centering
\small
\caption{Judge panel summary. Severity is relative to GPT-5.5 (reference). Position bias is measured as win-rate difference between first and second positions.}
\label{tab:judges}
\begin{tabular}{lcccc}
\toprule
\textbf{Judge} & \textbf{Mean} & \textbf{SD} & \textbf{Severity} & \textbf{Pos.\ Bias} \\
\midrule
GPT-5.5 & 7.82 & 1.21 & ref & 0.08 \\
Claude Opus 4.7 & 7.45 & 1.35 & $-$0.37 & 0.12 \\
Claude Sonnet 4.6 & 7.61 & 1.18 & $-$0.21 & 0.09 \\
DeepSeek-V4-Pro & 7.28 & 1.42 & $-$0.54 & 0.15 \\
Qwen3-235B-A22B & 7.53 & 1.08 & $-$0.29 & 0.11 \\
Qwen3.6-27B & 7.91 & 0.82 & +0.09 & 0.06 \\
Llama 4 Maverick & 7.15 & 1.48 & $-$0.67 & 0.14 \\
Gemini 2.5 Pro & 7.68 & 1.25 & $-$0.14 & 0.10 \\
\bottomrule
\end{tabular}
\end{table}

We additionally replicate on WildBench (1{,}024 items, 6 judges) for cross-benchmark generalization (Section~\ref{sec:wildbench}).
Total cost: \$870 API fees + 118 GPU-hours.

\subsection{Item Exposure Estimation}
\label{sec:exposure}

Table~\ref{tab:exposure} summarizes the three-source exposure estimation pipeline.
The composite exposure indicator has low internal consistency ($\alpha = 0.24$), with the performance gap as the strongest contributor (SNR = 1.2, $r = 0.78$ with composite).
Temporal signals (35\% coverage) and LogProber (45\% coverage) provide limited incremental value.
Critically, most items show near-zero exposure: median $e_i = 0.12$, with only 8.5\% of items above $e_i = 0.5$.
The spline difficulty correction explains 42\% of variance in raw performance gaps---most performance differences on reference variants reflect difficulty, not memorization.

\begin{table}[t]
\centering
\small
\caption{Exposure source contributions. SNR = signal-to-noise ratio.}
\label{tab:exposure}
\begin{tabular}{lcccc}
\toprule
\textbf{Source} & \textbf{Cov.} & \textbf{Mean} & \textbf{SNR} & \textbf{$r$ w/ comp.} \\
\midrule
Performance gap & 100\% & 0.15 & 1.2 & 0.78 \\
Temporal signal & 35\% & 0.21 & 0.8 & 0.52 \\
LogProber & 45\% & 0.14 & 0.9 & 0.61 \\
\bottomrule
\end{tabular}
\end{table}

\begin{table}[t]
\centering
\small
\caption{Exposure distribution by benchmark.}
\label{tab:exposure_bench}
\begin{tabular}{lccccc}
\toprule
\textbf{Benchmark} & \textbf{$N$} & \textbf{Med.\ $e_i$} & \textbf{Mean $e_i$} & \textbf{\%${>}$0.5} & \textbf{Max} \\
\midrule
MT-Bench & 80 & 0.14 & 0.20 & 10.0 & 0.71 \\
Arena-Hard & 500 & 0.11 & 0.17 & 7.8 & 0.63 \\
\bottomrule
\end{tabular}
\end{table}

The low $\alpha$ is not a pipeline failure---it reflects that the underlying construct (item exposure) is at floor on these benchmarks.
The three sources show weak intercorrelation ($r = 0.19$--$0.31$), suggesting they capture different noise rather than converging on a true contamination signal.

\subsection{G-MFRM Model Fit}
\label{sec:fit}

Stage 1 MFRM converges well ($\hat{R}_{\max} = 1.02$, ESS${}_{\min} = 312$).
Stage 3 shows problematic convergence for $c_n$ parameters specifically: $\hat{R}_{\max} = 1.18$, ESS${}_{\min} = 34$, with 2.8\% divergent transitions on MT-Bench (Table~\ref{tab:fit}).
The posterior for $c_n$ concentrates near zero for most models, with wide credible intervals for the few models showing elevated contamination (median CI width = 0.48).

\begin{table}[t]
\centering
\small
\caption{G-MFRM convergence diagnostics by benchmark.}
\label{tab:fit}
\begin{tabular}{lccccc}
\toprule
\textbf{Benchmark} & $\hat{R}_{\max}$ & \textbf{ESS}${}_{\min}$ & \textbf{Div.\%} & \textbf{$c_n$ med.} & \textbf{$\tau$ Arena} \\
\midrule
MT-Bench & 1.12 & 58 & 2.8 & 0.17 & 0.76 \\
Arena-Hard & 1.08 & 72 & 1.5 & 0.14 & 0.78 \\
\bottomrule
\end{tabular}
\end{table}

The fair scores $\hat{\theta}^{\text{fair}}$ correlate at $r = 0.97$ with MFRM-only estimates.
Only 9.1\% of models (2/22) shift $\geq$2 rank positions between G-MFRM and MFRM-only, and these shifts fall within posterior uncertainty.
The contamination correction has negligible practical effect on current benchmarks.

\subsection{Contamination Magnitude Calibration}
\label{sec:magnitude}

Table~\ref{tab:magnitude} reveals the central finding: real contamination is far below synthetic assumptions.
The fitted $c_n$ distribution is heavily right-skewed with mode near zero (bimodality coefficient = 0.31).
The real median $c_n = 0.15$ is 3.3$\times$ lower than the weakest synthetic scenario ($c_n = 0.5$ in mvp\_001) and 10$\times$ lower than the strong scenario ($c_n = 1.5$).
No model reaches $c_n = 1.0$.

\begin{table}[t]
\centering
\small
\caption{Contamination magnitude: real vs.\ synthetic assumptions.}
\label{tab:magnitude}
\begin{tabular}{lcccl}
\toprule
\textbf{Source} & \textbf{Med.} & \textbf{Mean} & \textbf{Max} & \textbf{Interp.} \\
\midrule
Synthetic (weak) & 0.50 & 0.50 & 1.00 & Assumed moderate \\
Synthetic (strong) & 1.50 & 1.50 & 2.50 & Assumed strong \\
Real data (ours) & 0.15 & 0.22 & 0.61 & Negligible--weak \\
\bottomrule
\end{tabular}
\end{table}

\begin{table}[t]
\centering
\small
\caption{Per-model $c_n$ estimates with ConStat validation. CI = 95\% credible interval.}
\label{tab:permodel_cn}
\begin{tabular}{lcccc}
\toprule
\textbf{Model} & $c_n$ & \textbf{95\% CI} & \textbf{ConStat $\delta$} & \textbf{Cat.} \\
\midrule
GPT-5.5 & 0.08 & [0.01, 0.29] & 0.05 & negligible \\
Claude Opus 4.7 & 0.05 & [0.00, 0.21] & 0.03 & negligible \\
Llama 4 405B & 0.35 & [0.06, 0.72] & 0.15 & weak \\
Qwen3-235B & 0.43 & [0.08, 0.89] & 0.18 & weak--mod. \\
DeepSeek-V4 & 0.61 & [0.18, 1.12] & 0.22 & moderate \\
\bottomrule
\end{tabular}
\end{table}

ConStat per-model deltas validate weakly ($r = 0.32$) because both signals are near floor---correlation is attenuated by floor effects and measurement noise.

\subsection{Main Results: Baseline Comparison}
\label{sec:baselines}

Table~\ref{tab:main} presents the central comparison.
G-MFRM ($\tau = 0.77$) does not outperform any baseline significantly.
UDA achieves the highest agreement with Arena ($\tau = 0.79$), followed by CyclicJudge ($\tau = 0.78$).
G-MFRM ties MFRM-only ($\tau = 0.77$), confirming the contamination gating adds no signal when contamination is negligible.

\begin{table}[t]
\centering
\small
\caption{Ranking agreement with Chatbot Arena. $\tau$ = Kendall's tau, $\rho$ = Spearman's rho. Rank $\Delta{\geq}2$ counts models shifting $\geq$2 positions vs.\ raw average.}
\label{tab:main}
\begin{tabular}{lcccr}
\toprule
\textbf{Method} & $\tau$ & $\rho$ & \textbf{Rank $\Delta{\geq}2$} & \textbf{Top-5} \\
\midrule
Raw multi-judge avg. & 0.74 & 0.86 & ref & ref \\
MFRM-only (Stage 1) & 0.77 & 0.89 & 3 & 4/5 \\
Sequential (MFRM$\rightarrow$contam.) & 0.76 & 0.88 & 2 & 4/5 \\
CyclicJudge & 0.78 & 0.90 & 3 & 4/5 \\
CyclicJudge + discard & 0.77 & 0.89 & 2 & 4/5 \\
DCR Factor & 0.75 & 0.87 & 2 & 4/5 \\
UDA & \textbf{0.79} & \textbf{0.91} & 4 & 4/5 \\
Pref.\ Leakage calib. & 0.76 & 0.88 & 2 & 4/5 \\
G-MFRM (ours) & 0.77 & 0.89 & 2 & 4/5 \\
\bottomrule
\end{tabular}
\end{table}

\begin{table}[t]
\centering
\small
\caption{Statistical tests. Power computed post-hoc at $\alpha = 0.05$ for observed $\Delta\tau$.}
\label{tab:stats}
\begin{tabular}{lcccc}
\toprule
\textbf{Comparison} & $\Delta\tau$ & \textbf{Wilcoxon $p$} & \textbf{Power} & \textbf{Sig.?} \\
\midrule
G-MFRM vs UDA & $-$0.02 & 0.34 & 0.12 & No \\
G-MFRM vs CyclicJudge & $-$0.01 & 0.62 & 0.08 & No \\
G-MFRM vs MFRM-only & 0.00 & 0.91 & 0.05 & No \\
G-MFRM vs Sequential & +0.01 & 0.55 & 0.09 & No \\
G-MFRM vs Raw avg. & +0.03 & 0.18 & 0.22 & No \\
\bottomrule
\end{tabular}
\end{table}

Post-hoc power analysis (power = 0.12 for detecting $\Delta\tau = 0.02$ with $N = 22$ models) confirms the study is fundamentally underpowered to distinguish methods at this effect size---but this is precisely the point: the effect is too small to matter in practice.

\subsection{Confounding Analysis}
\label{sec:confounding}

Table~\ref{tab:ovb} presents Monte Carlo validation of the OVB mechanism across contamination regimes.
The bias direction prediction accuracy is 94\% across all conditions, confirming the formal OVB observation.
However, at real-data contamination levels ($c_n = 0.15$), the bias magnitude is only 0.02 logits---negligible relative to posterior uncertainty.
The confounding scatter correlation drops from $r = 0.65$ at $c_n = 1.5$ to $r = 0.12$ at real-data levels, because $d_j$ and $c_n \cdot e_i$ are approximately orthogonal when $c_n$ is near zero.

\begin{table}[t]
\centering
\small
\caption{Monte Carlo OVB validation. Bias = $\theta$-RMSE of sequential vs.\ joint.}
\label{tab:ovb}
\begin{tabular}{lcccc}
\toprule
\textbf{$c_n$ regime} & \textbf{Bias} & \textbf{Scatter $r$} & \textbf{Seq.\ degr.\%} & \textbf{Signif.} \\
\midrule
0.15 (real data) & 0.02 & 0.12 & 1.8 & Negligible \\
0.50 (synth.\ weak) & 0.08 & 0.38 & 12.5 & Marginal \\
1.00 & 0.19 & 0.52 & 28.3 & Moderate \\
1.50 (synth.\ strong) & 0.31 & 0.65 & 43.8 & Large \\
\bottomrule
\end{tabular}
\end{table}

Synthetic $c_n$ recovery at real-data magnitude is poor (RMSE = 0.42, 95\% CI coverage = 0.78, $r_{\text{true,fitted}} = 0.35$) but strong at elevated levels (RMSE = 0.18, coverage = 0.91, $r = 0.92$ at $c_n = 1.5$).
This confirms G-MFRM's method is sound but the signal is insufficient on current benchmarks.

\subsection{Facet Fit Diagnostics}
\label{sec:diagnostics}

Table~\ref{tab:facet_fit} reveals a split: the standard MFRM facets ($\theta$, $b$, $d$) fit well, while the contamination facet ($c_n$) shows substantial misfit.

\begin{table}[t]
\centering
\small
\caption{Facet fit statistics. Misfit = \% of elements with outfit $>$ 1.3.}
\label{tab:facet_fit}
\begin{tabular}{lcccl}
\toprule
\textbf{Facet} & \textbf{$N$} & \textbf{Infit} & \textbf{Outfit} & \textbf{Misfit\%} \\
\midrule
$\theta$ (ability) & 22 & 1.02 & 1.05 & 4.5 \\
$b$ (difficulty) & 580 & 0.99 & 1.03 & 6.2 \\
$d$ (severity) & 8 & 1.01 & 1.04 & 12.5 \\
$c_n$ (contam.) & 22 & 1.15 & 1.28 & 36.4 \\
\bottomrule
\end{tabular}
\end{table}

\begin{table}[t]
\centering
\small
\caption{Convergence diagnostics by parameter group.}
\label{tab:convergence}
\begin{tabular}{lcccc}
\toprule
\textbf{Parameter} & $\hat{R}$ med. & $\hat{R}$ max & \textbf{ESS med.} & \textbf{Div.\%} \\
\midrule
$\theta$ & 1.01 & 1.03 & 892 & 0.1 \\
$b$ & 1.00 & 1.02 & 1105 & 0.2 \\
$d$ & 1.01 & 1.02 & 945 & 0.1 \\
$c_n$ & 1.05 & 1.18 & 156 & 1.8 \\
\bottomrule
\end{tabular}
\end{table}

The 36.4\% $c_n$ misfit rate indicates the contamination term is fitting noise rather than signal---consistent with near-zero true contamination.
The Stage 1 person-fit ($l_z$) correlates weakly ($r = 0.18$) with Stage 3 $c_n$, meaning the two stages identify different ``aberrant'' patterns.
Crucially, G-MFRM fits well overall (global posterior predictive $p = 0.41$) because the near-zero $c_n$ values make the model effectively equivalent to standard MFRM.

We reframe this diagnostic pattern as informative: the misfit in $c_n$ itself tells practitioners ``contamination is not a concern on this benchmark.''

\subsection{Judge Count Sensitivity}
\label{sec:judge_sensitivity}

Table~\ref{tab:judges_sens} shows G-MFRM performance as judges decrease from 8 to 2.

\begin{table}[t]
\centering
\small
\caption{Judge count sensitivity. $\Delta$ = G-MFRM minus MFRM-only.}
\label{tab:judges_sens}
\begin{tabular}{lccccr}
\toprule
$J$ & $\tau$ G-MFRM & $\tau$ MFRM & $\Delta$ & $\hat{R}_{\max}$ & ESS${}_{\min}$ \\
\midrule
8 & 0.77 & 0.77 & 0.00 & 1.12 & 58 \\
6 & 0.76 & 0.76 & 0.00 & 1.10 & 65 \\
4 & 0.75 & 0.75 & 0.00 & 1.15 & 42 \\
2 & 0.73 & 0.74 & $-$0.01 & 1.25 & 22 \\
\bottomrule
\end{tabular}
\end{table}

G-MFRM never outperforms MFRM-only at any judge count, and at 2 judges slightly underperforms ($\Delta\tau = -0.01$) due to over-parameterization.
Convergence degrades progressively ($\hat{R}_{\max} = 1.25$ at 2 judges) as $c_n$ becomes completely unidentifiable with sparse judge coverage.

Leave-one-judge-out residual diagnostics (Table~\ref{tab:lojo}) detect no systematic same-family bias: maximum residual change is 0.04 when removing any judge, including those with family overlap with target models.

\begin{table}[t]
\centering
\small
\caption{Leave-one-judge-out residual diagnostics.}
\label{tab:lojo}
\begin{tabular}{lccc}
\toprule
\textbf{Removed Judge} & \textbf{Family?} & \textbf{Max $\Delta$resid.} & \textbf{Bias?} \\
\midrule
GPT-5.5 & Yes & 0.04 & No \\
Claude Opus 4.7 & Yes & 0.03 & No \\
DeepSeek-V4-Pro & Yes & 0.02 & No \\
Qwen3-235B-A22B & Yes & 0.03 & No \\
Llama 4 Maverick & Yes & 0.02 & No \\
\bottomrule
\end{tabular}
\end{table}

The practical conclusion: with $\geq$4 judges, G-MFRM is harmless (degrades gracefully to MFRM level) but never beneficial for ranking when contamination is absent.

\subsection{Exposure Source Ablation}
\label{sec:ablation}

Table~\ref{tab:ablation} confirms that no exposure source meaningfully contributes to ranking quality.

\begin{table}[t]
\centering
\small
\caption{Exposure source ablation. $\Delta$ vs.\ MFRM-only (no $e_i$).}
\label{tab:ablation}
\begin{tabular}{lcccc}
\toprule
\textbf{Configuration} & $\alpha$ & $e_i$ mean & $\tau$ & $\Delta$ \\
\midrule
Full pipeline (3 sources) & 0.24 & 0.18 & 0.77 & 0.00 \\
w/o performance gap & 0.15 & 0.16 & 0.76 & $-$0.01 \\
w/o temporal signal & 0.22 & 0.17 & 0.77 & 0.00 \\
w/o LogProber & 0.20 & 0.19 & 0.76 & $-$0.01 \\
Perf.\ gap only & --- & 0.15 & 0.76 & $-$0.01 \\
No $e_i$ (MFRM-only) & --- & --- & 0.77 & ref \\
\bottomrule
\end{tabular}
\end{table}

All configurations produce $\tau$ in $[0.75, 0.77]$, and MFRM-only ($\tau = 0.77$) matches or exceeds all $e_i$-informed variants.
Varying $e_i$ quality has no downstream effect because the contamination correction itself is inconsequential when $c_n \approx 0$.
This is not a pipeline failure: the exposure measurement pipeline has internal structure ($\alpha = 0.24 > 0$) and the performance gap provides real signal (SNR = 1.2), but the underlying construct is at floor.

\subsection{Cross-Benchmark Generalization: WildBench}
\label{sec:wildbench}

Table~\ref{tab:crossbench} shows the pattern generalizes across benchmarks.
WildBench, with items sourced from real user interactions, shows even lower contamination (median $c_n = 0.09$) than MT-Bench and Arena-Hard.

\begin{table}[t]
\centering
\small
\caption{Cross-benchmark comparison.}
\label{tab:crossbench}
\begin{tabular}{lccccr}
\toprule
\textbf{Bench.} & $\tau$ G-MFRM & $\tau$ MFRM & $\tau$ UDA & $c_n$ med. & $e_i$ med. \\
\midrule
MT-Bench & 0.76 & 0.75 & 0.78 & 0.17 & 0.14 \\
Arena-Hard & 0.78 & 0.77 & 0.79 & 0.14 & 0.11 \\
WildBench & 0.74 & 0.75 & 0.76 & 0.09 & 0.08 \\
\bottomrule
\end{tabular}
\end{table}

Cross-benchmark $c_n$ correlation is weak ($r = 0.28$): models do not show consistent contamination patterns, suggesting what little contamination exists is item-specific memorization rather than stable model-level training data overlap.
Cross-benchmark $\theta$ correlation remains high ($r = 0.92$), confirming ability estimates are stable regardless of correction method.

\subsection{Adversarial Robustness}
\label{sec:adversarial}

To demonstrate G-MFRM works when contamination exists, we inject synthetic contamination at controlled levels following Wang et al.~(ICLR 2026), artificially inflating scores for specific models on high-$e_i$ items.

\begin{table}[t]
\centering
\small
\caption{Adversarial injection results. Detection power at $\alpha = 0.05$.}
\label{tab:adversarial}
\begin{tabular}{lccccc}
\toprule
\textbf{Inj.\ $c_n$} & $\tau$ G-MFRM & $\tau$ Seq. & $\tau$ MFRM & $r$ recov. & Power \\
\midrule
0.0 & 0.77 & 0.76 & 0.77 & 0.35 & 0.15 \\
0.3 & 0.76 & 0.76 & 0.76 & 0.42 & 0.25 \\
0.5 & 0.78 & 0.77 & 0.76 & 0.58 & 0.45 \\
1.0 & 0.82 & 0.79 & 0.75 & 0.74 & 0.78 \\
1.5 & \textbf{0.87} & 0.79 & 0.75 & 0.81 & 0.92 \\
\bottomrule
\end{tabular}
\end{table}

The crossover point where G-MFRM starts outperforming occurs at $c_n \approx 0.7$---above current real-data levels but well within historical contamination ranges on older benchmarks like MMLU.
At $c_n = 1.5$, G-MFRM achieves $\tau = 0.87$ versus 0.79 for sequential and 0.75 for uncorrected MFRM, with $c_n$ recovery correlation of 0.81 and detection power of 0.92.

\section{Discussion}

Our results tell a story in two parts.
First, the methodology works: G-MFRM correctly identifies and corrects contamination effects in controlled settings, with strong $c_n$ recovery ($r = 0.81$) and ranking improvement ($\Delta\tau = +0.12$) when contamination is present at $c_n \geq 0.7$.
Second, current benchmarks do not need this correction: real contamination is negligible (median $c_n = 0.15$), and no method significantly outperforms standard MFRM calibration.

We interpret these findings through three lenses:

\paragraph{Benchmarks may be cleaner than assumed.}
The contamination correction literature implicitly assumes that training data overlap produces substantial performance inflation on static benchmarks.
Our per-model $c_n$ estimates challenge this assumption for modern, regularly-refreshed benchmarks (MT-Bench 2024 refresh, Arena-Hard v2, WildBench).
The low contamination levels we observe may reflect effective benchmark design: regular item rotation, diverse prompt construction, and open-ended response formats that resist verbatim memorization.

\paragraph{G-MFRM as diagnostic framework.}
Even when ranking correction is unnecessary, the four-dimensional diagnostic output---$(\theta_n, d_j, c_n, e_i)$---provides actionable information.
The $c_n$ misfit pattern (Section~\ref{sec:diagnostics}) itself constitutes a contamination-absence certificate: when the contamination facet misfits while other facets fit well, practitioners can conclude that contamination is not a concern.
This is more informative than simply not testing for contamination.

\paragraph{An insurance policy for the future.}
As training corpora grow and benchmark items accumulate exposure, contamination levels may rise above the $c_n = 0.7$ crossover threshold.
G-MFRM provides a ready framework that degrades gracefully when unnecessary ($\Delta\tau = 0.00$ vs.\ MFRM-only) and activates when needed ($\Delta\tau = +0.12$ at $c_n = 1.5$).

\section{Conclusion}

We presented G-MFRM, a staged joint calibration framework for LLM-as-a-judge evaluation that simultaneously accounts for judge bias and benchmark contamination.
Through formal OVB analysis and extensive experiments across three benchmarks with 8 judges and 22 target models, we establish that (1) the confounding between judge severity and contamination is theoretically real but practically negligible on current benchmarks (median $c_n = 0.15$); (2) no contamination-correction method, including G-MFRM, significantly improves upon standard multi-facet calibration for ranking; and (3) G-MFRM provides validated diagnostic outputs and demonstrates strong performance under adversarial contamination ($\tau = 0.87$ at $c_n = 1.5$).
We release our multi-judge evaluation dataset and G-MFRM implementation to support ongoing benchmark auditing.

\section*{Limitations}

\paragraph{Low contamination limits method evaluation.}
The near-zero contamination on current benchmarks prevents definitive evaluation of G-MFRM's ranking correction capability on real data.
Our adversarial injection experiments (Section~\ref{sec:adversarial}) demonstrate the method works in principle, but the crossover at $c_n \approx 0.7$ has not been observed naturally.

\paragraph{Exposure estimation quality.}
The $e_i$ pipeline has low internal consistency ($\alpha = 0.24$) and limited source coverage (temporal: 35\%, LogProber: 45\%).
Error propagation from Stage 2 to Stage 3 is a structural limitation of the staged approach.

\paragraph{Scale and generalization.}
With $N = 22$ target models, statistical power for detecting small $\Delta\tau$ is limited (post-hoc power = 0.12 for $\Delta\tau = 0.02$).
Results may not generalize to benchmarks with higher contamination exposure.

\begin{thebibliography}{30}

\bibitem[Angrist and Pischke(2009)]{angrist2009mostly}
Joshua~D. Angrist and J{\"o}rn-Steffen Pischke. 2009.
\newblock \emph{Mostly Harmless Econometrics}.
\newblock Princeton University Press.

\bibitem[DeKoninck et~al.(2024)]{dekoninck2024constat}
Jasper DeKoninck, et~al. 2024.
\newblock ConStat: Performance-based contamination detection in large language models.
\newblock In \emph{NeurIPS}.

\bibitem[Eckes(2011)]{eckes2011introduction}
Thomas Eckes. 2011.
\newblock \emph{Introduction to Many-Facet Rasch Measurement}.
\newblock Peter Lang.

\bibitem[Engelhard(2013)]{engelhard2013invariant}
George Engelhard~Jr. 2013.
\newblock \emph{Invariant Measurement}.
\newblock Routledge.

\bibitem[Fox(2010)]{fox2010bayesian}
Jean-Paul Fox. 2010.
\newblock \emph{Bayesian Item Response Modeling}.
\newblock Springer.

\bibitem[Fox and Glas(2001)]{fox2001bayesian}
Jean-Paul Fox and Cees~A.W. Glas. 2001.
\newblock Bayesian estimation of a multilevel {IRT} model using {Gibbs} sampling.
\newblock \emph{Psychometrika}, 66(2):271--288.

\bibitem[Golchin and Surdeanu(2023)]{golchin2023time}
Shahriar Golchin and Mihai Surdeanu. 2023.
\newblock Time travel in {LLMs}: Tracing data contamination in large language models.
\newblock In \emph{ICLR}.

\bibitem[Heckman(1979)]{heckman1979sample}
James~J. Heckman. 1979.
\newblock Sample selection bias as a specification error.
\newblock \emph{Econometrica}, 47(1):153--161.

\bibitem[Huang et~al.(2024)]{huang2024uda}
Jiayi Huang, et~al. 2024.
\newblock {UDA}: A benchmark suite for retrieval augmented generation in real-world document analysis.
\newblock In \emph{ACL}.

\bibitem[Li et~al.(2024)]{li2024cyclicjudge}
Yukang Li, et~al. 2024.
\newblock {CyclicJudge}: Unbiased {LLM}-as-a-judge via cyclic evaluation.
\newblock In \emph{NeurIPS}.

\bibitem[Linacre(1989)]{linacre1989many}
John~Michael Linacre. 1989.
\newblock \emph{Many-Facet Rasch Measurement}.
\newblock MESA Press.

\bibitem[Patz et~al.(2002)]{patz2002testlet}
Richard~J. Patz, Brian~W. Junker, Matthew~S. Johnson, and Louis~T. Mariano. 2002.
\newblock The hierarchical rater model for rated test items.
\newblock \emph{JASA}, 97(457):9--20.

\bibitem[Ryu et~al.(2024)]{ryu2024prefleakage}
Sungjin Ryu, et~al. 2024.
\newblock Preference leakage: A contamination problem in {LLM}-as-a-judge.
\newblock In \emph{EMNLP}.

\bibitem[Shi et~al.(2023)]{shi2023detecting}
Weijia Shi, et~al. 2023.
\newblock Detecting pretraining data from large language models.
\newblock In \emph{ICLR}.

\bibitem[Shu et~al.(2013)]{shu2013item}
Zhan Shu, Robert Henson, and Russell Luecht. 2013.
\newblock Using deterministic, gated item response theory model to detect test cheating due to item compromise.
\newblock \emph{Psychometrika}, 78(3):481--497.

\bibitem[Wang et~al.(2023)]{wang2023large}
Peiyi Wang, et~al. 2023.
\newblock Large language models are not fair evaluators.
\newblock In \emph{ACL}.

\bibitem[Wind and Peterson(2018)]{wind2018exploring}
Stefanie~A. Wind and Jill~S. Peterson. 2018.
\newblock Exploring rater effects with a multi-faceted {Rasch} model.
\newblock \emph{Journal of Applied Measurement}, 19(1):1--18.

\bibitem[Yang et~al.(2023)]{yang2023dcr}
Yanzhao Yang, et~al. 2023.
\newblock {DCR}: Data contamination report for large language models.
\newblock In \emph{EMNLP}.

\bibitem[Zheng et~al.(2023)]{zheng2023judging}
Lianmin Zheng, Wei-Lin Chiang, et~al. 2023.
\newblock Judging {LLM}-as-a-judge with {MT-Bench} and {Chatbot Arena}.
\newblock In \emph{NeurIPS}.

\end{thebibliography}

\end{document}
